\documentclass[troisieme]{fiche}
\metadonnees{15}{Cours et va voir}{Bloc C — Décrire et comparer}

\begin{document}
\entete

\objectifs{%
\textbf{Langue} : bilan du groupe nominal ; thème.\par
\textbf{Texte} : Rudyard Kipling, \og Rikki-Tikki-Tavi \fg, \emph{The Jungle Book}
(1894).\par
\textbf{Durée} : 1 h — exercices 20 min, texte et questions 25 min, version 15 min.}

%% ===================================================================
\exercice[12 min]{Bilan du groupe nominal}

\rappel{\textbf{Rappel.} Tout ce qui décrit le nom se met devant lui, sans accord ; ce qui le
limite se met derrière.}

\consigne{Complétez d'après l'indication en italique.}

\begin{anglais}
\begin{enumerate}[label=\arabic*.,itemsep=2.5mm,leftmargin=*]
  \item \trou{Teddy's} father picked him up between his finger and thumb.
    \emph{(le père de Teddy)}
  \item It is \trou{the hardest} thing in the world to frighten a mongoose.
    \emph{(le plus difficile)}
  \item He was rather like \trou{a} little cat, but quite like \trou{a} weasel.
    \emph{(articles)}
  \item There wasn't \trou{any} water left in the ditch. \emph{(pas d'eau)}
  \item The boy \trou{who} found him was called Teddy. \emph{(relatif)}
  \item A mongoose has \trou{a longer} tail than a weasel. \emph{(plus longue)}
  \item \trou{Many} birds live in the garden, but there is not \trou{much} grass.
    \emph{(beaucoup)}
  \item They warmed him at \trou{the Englishman's} fire. \emph{(le feu de l'Anglais)}
  \item He had \trou{a few} friends in the garden: Darzee and Chuchundra.
    \emph{(quelques)}
  \item \trou{Mongooses} are curious animals. \emph{(les mangoustes en général)}
\end{enumerate}
\end{anglais}

\note[Le génitif]{\emph{'s} pour un possesseur animé (1, 8) ; \emph{of} pour une chose :
\emph{the end of the garden}. Au pluriel, l'apostrophe seule : \emph{the boys' room}.}

\note[Les articles]{\emph{A} pour une chose parmi d'autres (3, 6), \emph{the} pour celle dont
on a déjà parlé (4, 5). Pas d'article du tout pour une généralité au pluriel (10) :
\emph{mongooses are curious}, et non \emph{the mongooses are curious}, qui désignerait des
mangoustes précises.}

\note[Quantité]{\emph{Many} + pluriel, \emph{much} + indénombrable (7). \emph{A few} +
pluriel (9). \emph{Any} dans une phrase négative (4).}

%% ===================================================================
\exercice[8 min]{Thème}
\consigne{Traduisez en anglais.}

\begin{quote}\itshape
Ce soir-là, la mangouste voulut dormir dans le lit de Teddy. La mère du garçon n'était pas
contente : c'était un animal sauvage, disait-elle, et il pouvait le mordre. Mais le père
répondit qu'un enfant était plus en sécurité avec une mangouste qu'avec n'importe quel chien.
Rikki-tikki ne mordit personne. Il passa la nuit sur l'oreiller, et le lendemain matin il
sortit explorer le jardin, qui était le plus grand qu'il eût jamais vu.
\end{quote}

\lignes{9}

\corr{\textbf{Proposition.} That evening the mongoose wanted to sleep in Teddy's bed. The
boy's mother was not pleased: it was a wild animal, she said, and it could bite him. But the
father answered that a child was safer with a mongoose than with any dog. Rikki-tikki did not
bite anybody. He spent the night on the pillow, and the next morning he went out to explore
the garden, which was the biggest he had ever seen.}

\note[Quatre points]{\og le lit de Teddy \fg, \og la mère du garçon \fg{} : génitif, le
possesseur devant. \og il pouvait le mordre \fg{} : possibilité, \emph{could} ou \emph{might},
non \emph{was able to}. \og n'importe quel \fg{} : \emph{any} dans une phrase affirmative.
\og ne mordit personne \fg{} : une seule négation en anglais — \emph{did not bite anybody} ou
\emph{bit nobody}.}

\note[Le superlatif final]{\og le plus grand qu'il eût jamais vu \fg{} :
\emph{the biggest he had ever seen}. Pas de pronom relatif, et un pluperfect à cause de
\emph{ever}. La virgule devant \emph{which} est obligatoire : le jardin est déjà identifié.}

%% ===================================================================
\newpage
\section{Texte}

\noindent\small\textit{Dans un bungalow du nord de l'Inde. Rikki-tikki-tavi est une mangouste,
petit animal qui chasse les serpents.}\normalsize

\begin{texte}
He was a \gl[a mongoose]{mongoose}{une mangouste}, rather like a little cat in his fur and his
tail, but quite like a \gl[a weasel]{weasel}{une belette} in his head and his habits. His eyes
and the end of his \gl[restless]{restless}{sans repos, remuant} nose were pink. He could
scratch himself anywhere he pleased with any leg, front or back, that he chose to use. He
could \gl[to fluff up]{fluff up}{faire bouffer} his tail till it looked like a bottle brush,
and his war cry as he \gl[to scuttle]{scuttled}{détaler} through the long grass was:
``Rikk-tikk-tikki-tikki-tchk!''

One day, a high summer \gl[a flood]{flood}{une crue, une inondation} washed him out of the
\gl[a burrow]{burrow}{un terrier} where he lived with his father and mother, and carried him,
kicking and \gl[to cluck]{clucking}{glousser}, down a roadside
\gl[a ditch]{ditch}{un fossé}. He found a little \gl[a wisp]{wisp}{une touffe, un brin} of
grass floating there, and \gl[to cling, clung]{clung}{s'accrocher} to it till he lost his
senses. When he \gl[to revive]{revived}{revenir à soi} he was lying in the hot sun on the
middle of a garden path, very \gl[draggled]{draggled}{trempé et crotté} indeed, and a small
boy was saying, ``Here's a dead mongoose. Let's have a \gl[a funeral]{funeral}{un
enterrement}.''

``No,'' said his mother, ``let's take him in and dry him. Perhaps he isn't really dead.''

They took him into the house, and a big man picked him up between his finger and
\gl[a thumb]{thumb}{un pouce} and said he was not dead but half \gl[to choke]{choked}{étouffer}.
So they wrapped him in \gl[cotton wool]{cotton wool}{du coton hydrophile}, and warmed him over
a little fire, and he opened his eyes and \gl[to sneeze]{sneezed}{éternuer}.

``Now,'' said the big man (he was an Englishman who had just moved into the bungalow),
``don't frighten him, and we'll see what he'll do.''

It is the hardest thing in the world to frighten a mongoose, because he is eaten up from nose
to tail with curiosity. The \gl[a motto]{motto}{une devise} of all the mongoose family is
``Run and find out,'' and Rikki-tikki was a true mongoose. He looked at the cotton wool,
decided that it was not good to eat, ran all round the table, sat up and put his fur in order,
scratched himself, and jumped on the small boy's shoulder.

``Don't be frightened, Teddy,'' said his father. ``That's his way of making friends.''

``Ouch! He's \gl[to tickle]{tickling}{chatouiller} under my chin,'' said Teddy.

Rikki-tikki looked down between the boy's collar and neck, \gl[to snuff]{snuffed}{renifler} at
his ear, and climbed down to the floor, where he sat rubbing his nose.

``Good gracious,'' said Teddy's mother, ``and that's a wild creature! I suppose he's so
\gl[tame]{tame}{apprivoisé} because we've been kind to him.''

``All mongooses are like that,'' said her husband. ``If Teddy doesn't pick him up by the tail,
or try to put him in a \gl[a cage]{cage}{une cage}, he'll run in and out of the house all day
long. Let's give him something to eat.''
\end{texte}
\source{Rudyard Kipling, \og Rikki-Tikki-Tavi \fg, \emph{The Jungle Book}, Londres, Macmillan,
1894.}

\partiel[15 min]{Questions}
\consigne{Répondez aux deux premières questions en français, aux deux dernières en anglais.}

\begin{enumerate}[label=\arabic*.,itemsep=2.5mm,leftmargin=*]
  \item Comment Rikki-tikki se retrouve-t-il chez les Anglais ?
    \lignes{3}
    \corr{Une crue d'été l'a chassé du terrier où il vivait avec ses parents et l'a emporté
    dans un fossé. Il s'est accroché à une touffe d'herbe jusqu'à perdre connaissance, et
    s'est réveillé sur une allée de jardin, où un petit garçon l'a trouvé et où sa mère a
    demandé qu'on le rentre et qu'on le sèche.}
  \item Quel est le trait de caractère principal de la mangouste, d'après le texte ?
    \lignes{2}
    \corr{La curiosité, au point qu'il est très difficile de lui faire peur. La devise de
    toute la famille est \og Cours et va voir \fg.}
  \item \en{The motto of the mongoose family is ``Run and find out''. Find two moments in the
    text where Rikki-tikki obeys it.}
    \lignes{3}
    \corr{As soon as he wakes up he examines the cotton wool to see whether it can be eaten,
    runs all round the table and jumps on the boy's shoulder. He also looks down between
    Teddy's collar and neck and snuffs at his ear. He inspects everything before he fears
    anything.}
  \item \en{Find the sentence with a superlative and the relative clause in the second
    paragraph. Explain the construction of each.}
    \lignes{3}
    \corr{\emph{It is the hardest thing in the world to frighten a mongoose} : le vrai sujet
    est l'infinitif \emph{to frighten}, rejeté à la fin et annoncé par \emph{it} ; le
    complément du superlatif est introduit par \emph{in}. \emph{The burrow where he lived with
    his father and mother} : \emph{where} reprend un lieu — on pourrait dire \emph{in which},
    plus lourd.}
\end{enumerate}

%% ===================================================================
\section{Version}
\consigne{Traduisez le passage en français. Il suit le texte : Rikki-tikki vient de manger un
morceau de viande.}

\begin{version}
``There are more things to find out about in this house,'' he said to himself, ``than all my
family could find out in all their lives. I shall certainly stay and find out.''

He spent all that day \gl[to roam]{roaming}{parcourir, rôder dans} over the house. He nearly
\gl[to drown]{drowned}{noyer} himself in the bath-tubs, put his nose into the
\gl[ink]{ink}{l'encre} on a writing table, and burned it on the end of the big man's
\gl[a cigar]{cigar}{un cigare}, for he climbed up in the big man's \gl[a lap]{lap}{les genoux
(d'une personne assise)} to see how writing was done. At \gl[nightfall]{nightfall}{la tombée
de la nuit} he ran into Teddy's \gl[a nursery]{nursery}{une chambre d'enfant} to watch how
lamps were lighted, and when Teddy went to bed Rikki-tikki climbed up too. But he was a
restless companion, because he had to get up and attend to every noise all through the night,
and find out what made it.
\end{version}
\source{\emph{Ibid.}}

\lignes{14}

\corr{\textbf{Proposition de traduction.} \og Il y a plus de choses à découvrir dans cette
maison, se dit-il, que toute ma famille n'en découvrirait en toute une vie. Je vais
certainement rester, et voir. \fg{} Il passa toute cette journée-là à parcourir la maison. Il
faillit se noyer dans les baignoires, mit le nez dans l'encre d'une table à écrire et se le
brûla au bout du cigare du grand monsieur, car il lui avait grimpé sur les genoux pour voir
comment on écrivait. À la tombée de la nuit, il courut dans la chambre de Teddy pour regarder
comment on allumait les lampes, et quand Teddy alla se coucher, Rikki-tikki grimpa lui aussi.
Mais c'était un compagnon peu reposant, car il lui fallait se lever et s'occuper de chaque
bruit toute la nuit durant, et découvrir ce qui l'avait produit.}

\note[Choix de traduction 1]{\emph{to see how writing was done}, \emph{to watch how lamps were
lighted} : passifs sans agent. Le français passe par \og on \fg{} — \og comment on écrivait
\fg, \og comment on allumait les lampes \fg.}

\note[Choix de traduction 2]{\emph{he nearly drowned himself} : \emph{nearly} + prétérit =
\og faillit \fg. Traduire par \og presque \fg{} placerait l'adverbe au mauvais endroit.}

\note[Choix de traduction 3]{\emph{he had to get up} : \emph{must} n'ayant pas de passé, c'est
\emph{have to} qui porte l'obligation. Ici elle vient de son caractère, d'où \og il lui
fallait \fg.}

%% ===================================================================
\partiel{Lexique à mémoriser}
\begin{multicols}{2}\small
\begin{itemize}[label=--,itemsep=0.8mm,leftmargin=*]
  \item \textbf{a mongoose} — une mangouste \textit{(pl. mongooses)}
  \item \textbf{fur} — le pelage \textit{(fiche 10)}
  \item \textbf{a tail} — une queue
  \item \textbf{a burrow} — un terrier
  \item \textbf{a ditch} — un fossé
  \item \textbf{a flood} — une crue, une inondation
  \item \textbf{to frighten} — effrayer
  \item \textbf{curiosity} — la curiosité
  \item \textbf{to scratch} — gratter, griffer
  \item \textbf{tame} — apprivoisé ; \textit{contraire} wild
  \item \textbf{a shoulder} — une épaule
  \item \textbf{a noise} — un bruit
\end{itemize}
\end{multicols}

\end{document}
